\section{The Bipartite Double Cover and Strongly Leontovich Graphs}
\label{sec:doublecover}

A natural way to remove a loop from a target is to pass to the
bipartite double cover; Galvin et al.~\cite{galvin2026} note the
equivalent two-copy construction. We use this language to state the
preservation result and the even crossover threshold for the resulting
simple bipartite graph.

\begin{definition}[Bipartite Double Cover]
For any graph $H$, its bipartite double cover $B = H \times K_2$ is a
bipartite graph on $2|V(H)|$ vertices. The vertex set is $V(H) \times \{0, 1\}$,
with edges connecting $(u, 0)$ to $(v, 1)$ if and only if $(u, v) \in E(H)$.
A loop at vertex $u$ in $H$ corresponds to a standard edge between $(u, 0)$ and
$(u, 1)$ in $B$. Thus, $B$ is always a simple bipartite graph.
\end{definition}

\begin{theorem}
\label{thm:double_cover}
A graph $H$ is Leontovich if and only if its bipartite double cover $B = H
\times K_2$ is Leontovich. Furthermore, $H$ is Strongly Leontovich if and only
if $B$ is Strongly Leontovich.
\end{theorem}
\begin{proof}
Let $T$ be any connected tree. Because $T$ is bipartite, it admits exactly two
proper 2-colorings. Consequently, any homomorphism $\phi \colon T \to H$ lifts
to exactly two distinct homomorphisms $\Phi \colon T \to B$, corresponding to
the two choices of mapping the color classes of $T$ to the partitions $\{0, 1\}$
of $B$. Therefore, for any tree $T$, $\Hom(T, B) = 2 \Hom(T, H)$.

It follows immediately that for all $n$ and $d$:
\[
  \Hom(P_n, B) - \Hom(E_n^{(d)}, B) = 2 \Bigl[ \Hom(P_n, H) - \Hom(E_n^{(d)}, H) \Bigr]
\]
Because the homomorphism difference scales by a strict positive constant, the
crossover thresholds and the asymptotic ratio of the leading projection
coefficients $c_1(E_n)/c_1(P_n)$ are perfectly preserved.
\end{proof}

\subsection{A Simple Strongly Leontovich Graph}

The coefficient calculation in this section establishes that the looped symmetric
tree $\hat{T}(1, 35, 1, 50)$ on 1,822 vertices is Strongly Leontovich, with an asymptotic leading
coefficient ratio of

$\lim_{n \to \infty} \Hom(E_n^{(2)}, \hat{T}) / \Hom(P_n, \hat{T})
\approx 0.999953714414 < 1$.
Exact finite counts at $n=15001$ give the
still-unconverged ratio $0.999864892078$, reflecting the very slow decay of the
near-antipodal mode at eigenvalue $-7.282828$, whose magnitude is within a
factor $0.9997162$ of the Perron root $7.284896$.

By Theorem~\ref{thm:double_cover}, its bipartite double cover $B = \hat{T}
\times K_2$ inherits this exact property. $B$ is a simple, loopless, bipartite
graph on exactly 3,644 vertices. Because its asymptotic ratio is strictly less
than 1, $B$ is a \textbf{simple strongly Leontovich graph}, addressing the
simple-target existence question.

Furthermore, because $B$ is Strongly Leontovich, $\Hom(E_n, B) < \Hom(P_n, B)$
for all sufficiently large $n$, both even and odd. Exact integer computation of
the walk sequences reveals that $E_n$ overtakes $P_n$ on $B$ starting exactly at
the even threshold $n = 17{,}340$ and remains ahead through the verified range.
So even-$n$ crossovers do not require loops. Since $B$ is bipartite its
spectrum is exactly symmetric about zero, and that $\pm$ pairing is what
protects $P_n$ at even~$n$. What eventually defeats the protection is
near-degeneracy among the dominant magnitudes: the two largest positive
eigenvalues of $B$ are $7.284896$ and $7.282828$, a ratio of about
$0.99972$, so the competing mode decays slowly enough to survive until
$n = 17{,}340$. Equivalently, on the looped source $\hat{T}(1,35,1,50)$
the same fact appears as a near-antipodal eigenvalue $-7.282828$ just
inside the Perron root. This is not the second positive eigenvalue of
$\hat{T}$, $\lambda_2 \approx 5.842228$: at
$\lambda_2/\lambda_1 \approx 0.802$ that mode is negligible by
$n \approx 45$.

\paragraph{Disproof of the universal bipartite parity hypothesis.}
Appendix~\ref{app:sa} records a working hypothesis suggested by
extensive finite search: the 608,930,559 connected bipartite graphs on
$m \le 15$ vertices, together with all source trees through $n=22$
(including $5{,}623{,}756$ at $n=22$), showed no even-$n$ crossover. The
$n=17{,}340$ example above disproves the universal statement. The point
is scale, not a new mechanism: the relevant long-lived term is the
second dominant magnitude $7.282828$, not the smaller positive
eigenvalue $\lambda_2$.

\subsection{A Perturbed Non-Bipartite Certificate}

Adding a single edge inside one side of a bipartite strongly Leontovich
graph can produce a simple non-bipartite target. Here we use the looped
symmetric tree $\hat{T}(2, 34, 1, 48)$ on 3,403 vertices, whose
bipartite double cover $B' = \hat{T}(2, 34, 1, 48) \times K_2$ is a
simple bipartite strongly Leontovich graph on 6,806 vertices with
leading-coefficient ratio $0.999713000053\ldots$ and margin
$2.87 \times 10^{-4}$.

Adding a single edge inside one side of this bipartite graph creates an odd
cycle, so the perturbed target is strictly non-bipartite. We use
$\hat{T}(2, 34, 1, 48)$ rather than the smaller $\hat{T}(1,35,1,50)$ because
the larger asymptotic margin survives the perturbation more comfortably.
Adding an edge $e$ between two sibling terminal leaves of $B'$ gives the
6,806-vertex graph $B' + e$. The script
\texttt{scripts/verify\_perturbed\_cert.py} computes the coarsest equitable
partition by colour refinement, verifies equitability vertex-by-vertex with
exact integer arithmetic, and obtains an 18-cell quotient matrix. On this
quotient, the leading-coefficient ratio is
\[
  \lim_{n \to \infty}
  \frac{\Hom(E_n^{(2)}, B' + e)}{\Hom(P_n, B' + e)}
  \approx 0.999952140676 < 1.
\]
The same pipeline reproduces the unperturbed ratio
$0.999713000053\ldots$ for $B'$. A Rayleigh quotient estimate explains
why the perturbation is stable: the principal eigenvector mass on the
added leaves is tiny, and the measured shift is
$\Delta\lambda_1 \approx 1.77 \times 10^{-4}$.

To estimate the bound, \texttt{scripts/strong\_coeff.py}
(\texttt{certified\_leading\_ratio}) computes a Bauer--Fike residual
and a Davis--Kahan angle bound for the symmetrized 18-cell quotient,
both at 50 digits. Denote the resulting numerical estimate by
$\rho_{\mathrm{hi}}$. Numerically,
$\|r\|_2 \approx 7.5\times 10^{-50}$ against a spectral gap
$\lambda_1 - \lambda_2 \approx 5.1\times 10^{-3}$, propagating to
$\rho_{\mathrm{hi}} - \rho_2(B'+e) \approx 7.5\times 10^{-47}$ and
\[
  \rho_{\mathrm{hi}} \approx 0.99995214067566 < 1,
  \qquad 1 - \rho_{\mathrm{hi}} \approx 4.79\times 10^{-5}.
\]
The estimated error is roughly 41 orders of magnitude smaller than the
margin, so the numerical estimate $\rho_{\mathrm{hi}}$ stays below $1$. This is
high-precision numerical evidence (N) for a simple non-bipartite
strongly-Leontovich candidate in the leading-coefficient sense. It is a
residual-based estimate in that sense, not a directed-rounding
interval-arithmetic certificate.

% ════════════════════════════════════════════════════════════════════════
